\chapter{Autoencoder}
\label{ch14}

Autoencoders are among the most important deep-learning-based methods for anomaly detection. Unlike traditional supervised learning models that rely on labeled examples, autoencoders learn directly from the input data by attempting to reconstruct the original observations. Through this process, they learn \textbf{compact representations of normal patterns and can identify observations that differ substantially from those patterns.}

The key strength of an autoencoder lies in its ability to \textbf{compress high-dimensional data into a lower-dimensional latent representation} and then reconstruct the original input from that compressed form. This bottleneck structure forces the model to retain the most important information while discarding less relevant details. As a result, observations that do not conform to the learned structure often produce larger reconstruction errors, which can be used as anomaly scores.

Autoencoders have been successfully applied to a wide range of problems, including dimensionality reduction, image processing, denoising, feature extraction, and anomaly detection. Their ability to model complex nonlinear relationships makes them particularly valuable when traditional linear techniques such as Principal Component Analysis (PCA) are unable to capture important patterns in the data.

In this chapter, we introduce the architecture and intuition of autoencoders, including the encoder, decoder, and latent space. We explain how autoencoders compress and reconstruct data, compare them with PCA, and demonstrate why reconstruction error provides a natural mechanism for anomaly detection. We then develop a complete anomaly detection workflow using PyOD, including model development, threshold selection, result interpretation, and model evaluation. Finally, we discuss model aggregation techniques that can improve the stability and robustness of autoencoder-based anomaly detection systems.

\section{Learning Objectives}

After completing this chapter, you will be able to:

\begin{itemize}
    \item Explain the architecture and operation of an autoencoder, including the roles of the encoder, decoder, and latent (bottleneck) layer.
    \item Describe how autoencoders perform dimensionality reduction by learning compressed representations of data.
    \item Build and train an autoencoder model using PyOD.
Interpret anomaly detection results through descriptive statistics, score distributions, and confusion matrices.
    \item Evaluate the performance of an autoencoder-based anomaly detection model.
    \item Apply model aggregation techniques to improve model stability.
\end{itemize}

\section{Software Requirements}

Autoencoders are neural-network-based algorithms and require TensorFlow. To prevent unintended consequences, PyOD does not install TensorFlow automatically. This is explained on PyOD's installation page.
You need to install the following libraries. See the ``Software Requirements'' section in Chapter~1.

\begin{lstlisting}[language=Python]
!pip install tensorflow
!pip install pyod
!pip install combo
\end{lstlisting}


\section{Understanding an Autoencoder}

An autoencoder is a specialized type of neural network designed to reproduce its input as its output. In Figure \ref{fig:autoencoder}, the same image — such as the Mona Lisa — is used as both the input and the output. Thus, unlike most neural networks, which are typically used in supervised learning, autoencoders operate in an \textbf{unsupervised setting} because they do not require external target labels. The input itself serves as the target. For example.


\begin{figure}[htbp]
\centering
\includegraphics[width=.8\textwidth]{images/ch11_fig11.png}
\caption{Autoencoder architecture}
\label{fig:autoencoder}
\end{figure}


At first glance, it may seem unnecessary to train a model whose goal is simply to replicate its input. However, the key insight of an autoencoder lies not in the output layer, but in the internal representation learned in the hidden layers.

An effective autoencoder is typically designed with one or more hidden layers that contain fewer neurons than the input layer. This architectural constraint forces the network to \textbf{compress the input data into a lower-dimensional representation}, often referred to as the latent space or bottleneck layer. During this compression process, the model must retain only the most important features of the data while discarding irrelevant variations or noise.

If the hidden layers were larger than the input layer, the model could simply learn to copy the input directly, including any noise, without capturing meaningful structure. By contrast, a bottleneck architecture encourages the model to focus on the underlying patterns that best represent the data.

This ability to learn compact and informative representations makes autoencoders particularly useful for tasks such as dimensionality reduction, feature extraction, and anomaly detection.

Figure \ref{fig:autoencoder} also illustrates both the \textbf{encoding} and \textbf{decoding} processes of an autoencoder. The \textbf{encoding process} compresses the input data into a compact representation at the core (latent) layer. This process can be visualized as a funnel that is wide on the input side, where the full set of features enters, and narrows toward the center, where a reduced-dimensional representation is formed.

In contrast, the \textbf{decoding process} reconstructs the compressed representation to produce the output. This stage resembles an inverted funnel: narrow at the input of the decoder (the latent representation) and expanding outward as the data are reconstructed. The information contained in the original image is first compressed in the core hidden layer and then reconstructed as faithfully as possible. Consequently, the latent representation must preserve the essential structure and features of the input data.

In practice, the decoder is typically designed to mirror the encoder. This means that the number of layers and neurons in the decoding stage are structured symmetrically with respect to those in the encoding stage. Such symmetry is commonly adopted because it encourages efficient reconstruction while maintaining a balanced architecture.

The specific design choices for the encoder–decoder structure, including the number of layers and neurons, will be discussed in detail in the following sections.

\section{Applications of Autoencoders}

One of the earliest and most important applications of autoencoders is \textbf{dimensionality reduction}. A milestone study by Hinton and Salakhutdinov (2006) [1] demonstrated that autoencoders can achieve lower reconstruction error than the first 30 principal components obtained from Principal Component Analysis (PCA), while also improving the separation of data clusters. This result highlighted the ability of autoencoders to capture nonlinear structures in data more effectively than linear methods such as PCA.

Beyond dimensionality reduction, autoencoders have been widely applied in \textbf{computer vision} and \textbf{image-processing tasks}.
One notable example is \textbf{image colorization}, where autoencoders are trained to convert grayscale images into color images. As illustrated in Figure~\ref{fig:image_coloring}, the input to the model is a black-and-white image, while the target output is its corresponding color version.

\begin{figure}[htbp]
\centering
\includegraphics[width=.9\textwidth]{images/ch11_fig12.png}
\caption{An autoencoder model for image coloring}
\label{fig:image_coloring}
\end{figure}


During training, the model learns from a large set of paired examples consisting of grayscale images and their corresponding color images. Through this process, the autoencoder captures the underlying patterns and relationships needed to infer color information. Once trained, the model can generalize to new inputs and generate plausible color versions of previously unseen grayscale images.


Similarly, autoencoders can be applied to \textbf{image denoising} tasks. As illustrated in Figure~\ref{fig:image_denoising}, the input to the model is a noisy or blurry image, while the target output is a clean version of that image.
The model is trained on a large set of paired examples consisting of noisy images and their corresponding clean versions. 

Through this training process, the autoencoder learns to reconstruct the underlying structure of the image while filtering out noise.
Once trained, the model can generalize to new inputs and transform previously unseen noisy or blurry images into clearer versions.

Although more advanced models have since been developed for noise reduction, autoencoders have played a foundational role in this area and remain an important method for understanding and implementing image denoising techniques.

\begin{figure}[htbp]
\centering
\includegraphics[width=.8\textwidth]{images/ch11_fig13.png}
\caption{An autoencoder model for noise reduction}
\label{fig:image_denoising}
\end{figure}

\section{Comparing Autoencoders with PCA for Dimensionality Reduction}

When discussing dimensionality reduction, \textbf{Principal Component Analysis (PCA)} is often the most commonly referenced technique. While both PCA and autoencoders aim to reduce dimensionality, they differ significantly in their assumptions, methodology, and capabilities.

There are many effective methods for detecting outliers, including PCA. However, autoencoders offer several distinct advantages that make them particularly valuable for certain applications.

PCA is a linear technique that uses linear algebra to transform data into principal components, which are linear combinations of the original features. While PCA is effective for capturing dominant linear patterns in the data and reducing dimensionality, it is inherently limited to linear transformations and may struggle to represent more complex nonlinear structures.

In contrast, autoencoders employ \textbf{nonlinear activation functions and multilayer architectures} to learn nonlinear transformations of the input data. This capability enables autoencoders to capture intricate patterns and relationships that may not be accessible through linear methods such as PCA.
By stacking multiple layers, autoencoders can learn hierarchical representations and uncover complex feature interactions.

Furthermore, while PCA applies a single global transformation, autoencoders learn representations through a sequence of transformations across layers. This layered structure provides greater flexibility and adaptability, allowing the model to better accommodate complex data distributions.

As a result, autoencoders are particularly well suited for problems involving \textbf{high-dimensional and nonlinear data}, where capturing subtle patterns is essential for accurate anomaly detection.


\section{Why Autoencoders Can Identify Anomalies}

The primary goal of autoencoders is to learn an efficient compressed representation of the data and then reconstruct the original input from that representation. The ability to detect anomalies emerges naturally as a \textbf{byproduct} of this learning process.

The key lies in the bottleneck architecture of the autoencoder. During training, the encoder compresses the input into a lower-dimensional latent representation. Because the latent layer contains fewer dimensions than the original data, the network cannot simply memorize every detail of every observation. Instead, it must learn the most important and recurring patterns that characterize the underlying data distribution.

\textbf{Anomalies do not conform to the dominant patterns learned during training}. They cannot be encoded as efficiently within the latent space. Important information is lost during compression because the latent representation has been optimized to capture normal patterns rather than unusual ones. Consequently, when the decoder attempts to reconstruct the anomalous observation, the reconstruction is often inaccurate, resulting in a substantially larger reconstruction error.

\begin{figure}[htbp]
\centering
\includegraphics[width=.9\textwidth]{images/ch11_fig14.png}
\caption{Autoencoder reconstruction of normal and anomalous images}
\label{fig:autoencoder_anomaly}
\end{figure}


This behavior can be understood through a simple analogy shown in Figure~\ref{fig:autoencoder_anomaly}. Imagine training an autoencoder exclusively on \textbf{images of cats}. During training, the network learns features such as ear shapes, facial symmetry, whiskers, and fur textures. If another cat image is presented, the model can compress and reconstruct it effectively because it aligns with previously learned patterns. 

However, if a dog image is supplied instead, the latent representation does not contain an efficient description of dog-specific characteristics. The resulting reconstruction may resemble a distorted version of the input, leading to a noticeably higher reconstruction error.

Mathematically, the reconstruction error is often measured as the difference between the original input $x$ and the reconstructed output $\hat{x}$:

\[
\text{Reconstruction Error}
=
\|x-\hat{x}\|^2
\]

A low reconstruction error suggests that the observation is consistent with the patterns learned during training, whereas  \textbf{a high reconstruction error indicates that the observation may be an outlier or anomaly}.

Next, we apply autoencoders to anomaly detection.

\section{The Modeling Procedure}

We adopt a structured modeling procedure for the development, evaluation, and interpretation of results. This procedure consists of the following steps as shown in Figure \ref{fig:modeling_procedure}:

\begin{enumerate}
    \item \textbf{Model Development}
    
    Design and train the model to learn underlying patterns in the data and generate anomaly scores.
    
    \item \textbf{Threshold Determination}
    
    Once anomaly scores are obtained, a threshold must be selected to distinguish between normal and abnormal observations. For example, if prior knowledge suggests that anomalies should constitute no more than 1\% of the dataset, the threshold can be chosen such that approximately 1\% of the observations are classified as anomalies.
    
    \item \textbf{Descriptive Statistics of Normal and Abnormal Groups}
    
    Analyzing descriptive statistics for both groups is essential for evaluating model performance. These statistics help verify whether the model's outputs are consistent with domain expectations. For instance, if a feature is expected to have a higher mean in the abnormal group than in the normal group, but the results indicate otherwise, this discrepancy should be investigated. In such cases, the feature may need to be re-examined, transformed, or removed, and the model retrained.
\end{enumerate}

\begin{figure}[htbp]
\centering
\includegraphics[width=.8\textwidth]{images/ch11_fig14a.png}
\caption{Modeling procedure}
\label{fig:modeling_procedure}
\end{figure}

In an autoencoder, it is important that the hidden layers contain fewer neurons than the input layer. This architectural constraint encourages the model to learn a compact representation of the data.

\textbf{To illustrate this, we generate a synthetic dataset with up to 25 variables, which correspond to 25 neurons in the input layer of the autoencoder.} With this setup, we can experiment with different configurations of hidden layers and neuron sizes, allowing us to examine how the degree of compression affects the model's ability to capture essential patterns in the data.

\begin{lstlisting}[language=Python]
import numpy as np
import pandas as pd
from pyod.utils.data import generate_data
X_train, X_test, y_train, y_test = DGP(n_features=25)
\end{lstlisting}

\subsection{Step 1 --- Build the Model}

Here, we define a straightforward autoencoder with two hidden layers, each containing two neurons. This configuration is specified as \texttt{hidden\_neuron\_list=[2,2]}.

\begin{lstlisting}[language=Python]
from pyod.models.auto_encoder import AutoEncoder
atcdr = AutoEncoder(contamination=0.05, hidden_neuron_list =[2, 2])
atcdr.fit(X_train)
\end{lstlisting}

We then obtain the model predictions:

\begin{lstlisting}[language=Python]
results = evaluate_model(atcdr, X_train, X_test)
y_train_scores = results["train_scores"]
y_train_pred = results["train_pred"]

y_test_scores = results["test_scores"]
y_test_pred = results["test_pred"]
threshold = results["threshold"]
\end{lstlisting}

The next step is to determine a suitable threshold.

\subsection{Step 2 --- Determine a Reasonable Threshold}

By default, if the contamination rate is not specified, PyOD assumes a value of 0.10, indicating that approximately 10\% of the data are treated as anomalies.

In this example, the contamination rate is explicitly set to 0.05. Consequently, the threshold is computed such that the top 5\% of observations with the highest anomaly scores are classified as outliers.

For instance, the threshold corresponding to a 5\% contamination rate is 4.420. Observations with anomaly scores greater than this threshold are labeled as outliers, while those below the threshold are considered normal.

\begin{lstlisting}[language=Python]
print("The threshold for the defined contamination rate:" , threshold)
print("The training data:", count_stat(y_train_pred))
print("The test data:", count_stat(y_test_pred))
\end{lstlisting}

\begin{lstlisting}[language=Python]
The threshold for the defined contamination rate: 4.420323371886652
The training data: {0: 475,; 1: 25}
The training data: {0: 475,; 1: 25}
\end{lstlisting}


As discussed in previous chapters, the exact proportion of outliers in a dataset is often unknown. In such cases, an effective approach is to examine the distribution of anomaly scores to determine an appropriate threshold.

One common method is to use a histogram of the outlier scores, which provides a visual representation of their distribution. By identifying natural breaks or sharp declines in frequency, a reasonable cutoff point can be selected.

We now generate a histogram of the outlier scores.

\begin{lstlisting}[language=Python]
plot_hist(y_train_scores)
\end{lstlisting}

Figure~\ref{fig:autoencoder_histogram} presents a histogram of the outlier scores. By examining this distribution, we can identify a natural cutoff point where a clear separation emerges between the majority of normal observations and potential outliers. In this case, the histogram suggests a threshold of approximately 4.5.

This approach provides a practical method for threshold selection when the true proportion of outliers is unknown. By aligning the threshold with the underlying distribution of anomaly scores, we can make informed and data-driven decisions about which observations should be classified as anomalies.

\begin{figure}[htbp]
\centering
\includegraphics[width=.6\textwidth]{images/ch11_fig15.png}
\caption{The histogram of Autoencoder outlier scores}
\label{fig:autoencoder_histogram}
\end{figure}

\subsection{Step 3 --- Profile the Normal and Abnormal Groups}

Profiling the normal and outlier groups is a critical step in validating the effectiveness of an anomaly-detection model. This process involves comparing the feature statistics of these groups to ensure consistency with domain knowledge and prior expectations.

If domain knowledge suggests that the mean value of a particular feature should be higher in the outlier group than in the normal group, but the results indicate otherwise, this discrepancy should be carefully investigated. Potential actions include re-examining the feature, applying transformations, or removing the feature if it does not contribute meaningful information.

This iterative refinement process helps ensure that the model is both accurate and interpretable.

At the same time, the profiling process can uncover new patterns or insights within the data. When unexpected results arise, they may suggest that existing assumptions or domain knowledge should be reconsidered. In such cases, incorporating these new insights into the modeling process can deepen understanding and lead to improved model performance.

Thus, profiling serves not only as a validation mechanism but also as a valuable tool for knowledge discovery and model refinement.

\begin{lstlisting}[language=Python]
feature_list = ['F01','F02','F03','F04','F05','F06','F07','F08','F09','F10','F11','F12',
'F13','F14','F15','F16','F17','F18','F19','F20','F21','F22','F23','F24','F25']
descriptive_stat_threshold(X_train,y_train_scores, threshold, feature_list)
\end{lstlisting}

\begin{table}[ht]
\centering
\caption{Summary Statistics for Normal and Outlier Groups}
\label{tab:summary_stats}
\begin{tabular}{lrrrrrrrrrrr}
\hline
Group & Count & Count (\%) & F1 & F2 & F3 &
F4 & $\cdots$  & F23 &
F24 & F25 & Anomaly\_Score \\
\hline
Normal  & 475 & 95.0 & 2.99 & 2.99 & 2.98 & 3.00 &
$\cdots$ & 3.00 & 2.99 & 3.01 & 2.59 \\
Outlier & 25  & 5.0  & 5.02 & 4.79 & 5.05 & 4.87 &
$\cdots$  & 4.80 & 4.91 & 5.10 & 19.83 \\
\hline
\end{tabular}
\end{table}



Table~\ref{tab:summary_stats} summarizes the characteristics of the normal and abnormal groups, including both the counts and proportions of observations in each group. For clarity, all features should be clearly labeled with their respective names.
Several key insights can be drawn.

\begin{itemize}
\item \textbf{Size of the Outlier Group}

Once the threshold is established, the size of the outlier group is determined accordingly. This proportion provides a useful starting point for understanding the prevalence of anomalies within the dataset.

\item \textbf{Feature Statistics for Each Group}

The mean values of the features in the normal and outlier groups should be consistent with domain knowledge. For example, if prior knowledge suggests that certain features should have higher mean values in the outlier group, the results should reflect this expectation.

In this case, the means of the features in the outlier group are higher than those in the normal group, which should be validated against domain-specific insights.

\item \textbf{Average Anomaly Score}

The average anomaly score for the outlier group should be higher than that of the normal group, as anomalous observations typically receive higher scores.

Although the specific values may not require detailed interpretation, confirming this separation is essential for validating the model's effectiveness.

\end{itemize}

Because the ground truth is available, we can construct a confusion matrix to evaluate model performance. The results indicate that the model performs well, successfully identifying all 25 outliers.
\begin{lstlisting}[language=Python]
confusion_matrix(y_test,y_test_scores,threshold)
\end{lstlisting}

\begin{table}[ht]
\centering
\caption{Confusion Matrix}
\begin{tabular}{lcc}
\hline
 & \multicolumn{2}{c}{Pred} \\
\cline{2-3}
Actual & 0 & 1 \\
\hline
0 & 475 & 0 \\
1 & 0 & 25 \\
\hline
\end{tabular}
\label{tab:confusion_matrix}
\end{table}



\section{Aggregating to Achieve Model Stability}

As discussed by Aggarwal [2], one significant challenge in neural networks is their sensitivity to noise, which can lead to overfitting. To address overfitting and instability in model predictions, a common strategy is to train multiple models and aggregate their scores. During the aggregation process, Steps~2 and~3 of the modeling procedure should still be followed as previously described.

There are several methods for aggregating the outcomes of multiple models:

\begin{itemize}
    \item \textbf{Average:} Calculate the average of the scores from all detectors.
    
    \item \textbf{Maximum of Maximum (MOM):} Take the maximum score from each detector and then find the maximum value among those scores.
    
    \item \textbf{Average of Maximum (AOM):} First, find the maximum score from each detector and then average these maximum scores.
    
    \item \textbf{Maximum of Average (MOA):} Average the scores from all detectors and then take the maximum of these average scores.
\end{itemize}

You only need to choose one aggregation approach for analysis. In this discussion, we use the averaging method.

\begin{lstlisting}[language=Python]
from pyod.models.combination import (aom, moa, average, maximization)
from pyod.utils.utility import standardizer
from pyod.models.auto_encoder import AutoEncoder

atcdr1 = AutoEncoder(contamination=0.05,hidden_neuron_list=[2, 2])
atcdr2 = AutoEncoder(contamination=0.05,hidden_neuron_list=[10, 2, 10])
atcdr3 = AutoEncoder(contamination=0.05,hidden_neuron_list=[15, 10, 2, 10, 15])

# Prepare data frames to store results
train_scores = np.zeros([X_train.shape[0], 3])
test_scores = np.zeros([X_test.shape[0], 3])

atcdr1.fit(X_train)
atcdr2.fit(X_train)
atcdr3.fit(X_train)

# Store results
train_scores[:, 0] = atcdr1.decision_function(X_train)
train_scores[:, 1] = atcdr2.decision_function(X_train)
train_scores[:, 2] = atcdr3.decision_function(X_train)

test_scores[:, 0] = atcdr1.decision_function(X_test)
test_scores[:, 1] = atcdr2.decision_function(X_test)
test_scores[:, 2] = atcdr3.decision_function(X_test)

# Normalize scores
train_scores_norm, test_scores_norm = standardizer(train_scores,test_scores)

# Combination by average
y_train_by_average = average(train_scores_norm)
y_test_by_average = average(test_scores_norm)
\end{lstlisting}

To illustrate, we define three different models:

\begin{itemize}
    \item \textbf{Model \texttt{atcdr1}:}
    Two hidden layers, each containing two neurons.
    
    \item \textbf{Model \texttt{atcdr2}:}
    Three hidden layers with 10, 2, and 10 neurons, respectively.
    
    \item \textbf{Model \texttt{atcdr3}:}
    Five hidden layers with 15, 10, 2, 10, and 15 neurons, respectively.
\end{itemize}

In line with the convention of symmetric autoencoder architectures, Model \texttt{atcdr3} clearly exhibits symmetry. This arrangement helps ensure a balanced and effective encoding--decoding process.

The code performs the following steps:

\begin{enumerate}
    \item Prepare data frames to store model predictions.
    
    \item Train the three autoencoder models independently.
    
    \item Store the anomaly scores from each model.
    
    \item Standardize the predictions so they are on a comparable scale.
    
    \item Compute the average anomaly score across all models.
\end{enumerate}

The code below plots a histogram of the aggregated predictions.

\begin{lstlisting}[language=Python]
import matplotlib.pyplot as plt
plt.figure( figsize=(6, 4), dpi=80)
plt.hist( y_train_by_average, bins='auto')
plt.title("Combination by average")
plt.show()
\end{lstlisting}

The histogram shown in Figure~\ref{fig:autoencoder_ensemble_hist} indicates that a threshold of 0.0 is appropriate.

\begin{figure}[htbp]
\centering
\includegraphics[width=.6\textwidth]{images/ch11_fig16.png}
\caption{The histogram of the average prediction of the Autoencoder models}
\label{fig:autoencoder_ensemble_hist}
\end{figure}

We can derive descriptive statistics using the chosen threshold.

\begin{lstlisting}[language=Python]
descriptive_stat_threshold(X_train, y_train_by_average, 0.0)
\end{lstlisting}



\begin{table}[ht]
\centering
\caption{Summary of Normal and Outlier Groups}
\label{tab:avg_summary_stats}
\begin{tabular}{lrrrrrrrrrrr}
\hline
Group & Count & Count (\%) & F1 & F2 & F3 &
$\cdots$ & F22 & F23 & F24 &
F25 & Anomaly\_Score \\
\hline
Normal  & 471 & 94.2 & 2.99 & 2.99 & 2.97 &
$\cdots$ & 3.00 & 3.01 & 2.99 & 3.01 & -0.23 \\
Outlier & 29  & 5.8  & 4.74 & 4.49 & 4.79 &
$\cdots$ & 4.93 & 4.50 & 4.63 & 4.85 & 3.72 \\
\hline
\end{tabular}
\end{table}




Table~\ref{tab:avg_summary_stats} shows the results. The analysis identifies 25 observations as outliers. The interpretation is similar to that of Table~\ref{tab:summary_stats}.


Up to this point, the autoencoder modeling process for anomaly detection has been completed. We now proceed to examine the key hyperparameters used in neural networks.

\section{Neural Network Hyperparameter Tuning}

Let us print out the specifications of the model \texttt{atcdr}.



You can format it nicely in LaTeX using either a description list or a table. For a thesis/dissertation, I would recommend a table:

\begin{table}[htbp]
\centering
\caption{Autoencoder Hyperparameters}
\begin{tabular}{ll}
\hline
\textbf{Parameter} & \textbf{Value} \\
\hline
Batch Size & 32 \\
Contamination & 0.05 \\
Dropout Rate & 0.2 \\
Epochs & 100 \\
Hidden Activation & ReLU \\
Hidden Neurons & [2, 2] \\
L2 Regularizer & 0.1 \\
Loss Function & Mean Squared Error (MSE) \\
Optimizer & Adam \\
Output Activation & Sigmoid \\
Preprocessing & True \\
Random State & None \\
Validation Size & 0.1 \\
Verbose & 1 \\
\hline
\end{tabular}
\label{tab:autoencoder_hyperparameters}
\end{table}


If you want to display it exactly as a Python dictionary, use:

\begin{verbatim}
{
 'batch_size': 32,
 'contamination': 0.05,
 'dropout_rate': 0.2,
 'epochs': 100,
 'hidden_activation': 'relu',
 'hidden_neurons': [2, 2],
 'l2_regularizer': 0.1,
 'loss': mean_squared_error,
 'optimizer': 'adam',
 'output_activation': 'sigmoid',
 'preprocessing': True,
 'random_state': None,
 'validation_size': 0.1,
 'verbose': 1
}
\end{verbatim}



\section{Summary}

This chapter introduced autoencoders as a powerful deep learning approach for anomaly detection that learns normal data patterns through reconstruction rather than relying on labeled observations. By compressing high-dimensional data into a lower-dimensional latent representation and reconstructing the original inputs, autoencoders capture the most important structural characteristics of the data while filtering out less relevant information. We examined the fundamental components of an autoencoder, including the encoder, latent (bottleneck) layer, and decoder, and explored how reconstruction error can be used to identify anomalous observations. We demonstrated a complete anomaly detection workflow using PyOD, covering model development, training, threshold selection, result interpretation, performance evaluation, and model aggregation techniques to enhance the stability and robustness of autoencoder-based anomaly detection systems.

\section{Critical Thinking Questions}

\begin{enumerate}
\item Why does the bottleneck (latent) layer play a critical role in anomaly detection?
When would an autoencoder be preferred over PCA for dimensionality reduction and anomaly detection?
\item How should an anomaly detection threshold be selected when the true proportion of anomalies is unknown? Compare the use of a predefined contamination rate with the use of anomaly score distributions (such as histograms). What are the risks and benefits of each approach, and how might poor threshold selection affect business decision-making?
\item Why can combining multiple autoencoders improve anomaly detection performance and stability?

\end{enumerate}


\vspace{0.5cm}
\noindent \large{\textbf{References}}
\normalsize

\begin{enumerate}

\item
Hinton, G.~E.,
and Salakhutdinov, R.~R.
\newblock Reducing the Dimensionality of Data with Neural Networks.
\newblock \emph{Science},
313, 504--507, 2006.

\item
Aggarwal, C.~C.
\newblock \emph{Outlier Analysis}.
\newblock Springer, 2016.
\newblock ISBN: 978--3319475776.

\end{enumerate}



